\section{Sesgos y fairness}

En Doctor Maíz, fairness no se refiere a atributos demográficos. Se pregunta si
el clasificador mantiene un rendimiento parecido entre clases y condiciones de
captura agrícola. El reporte se generó únicamente sobre el holdout final y
calculó precision, recall y F1 por clase, entorno, fuente, tercil de resolución
y tercil de blur. Para evitar grupos diminutos se exigieron al menos 30 imágenes
por fuente y 20 en las demás dimensiones.

\begin{table}[H]
  \centering
  \caption{Diferencia entre el mayor y menor valor observado dentro de cada
  dimensión. Los gaps son descriptivos y no llevan signo.}
  \label{tab:fairness-gaps}
  \input{tables/fairness_gaps}
\end{table}

La brecha mayor apareció entre clases: 0.2584 en precision, 0.3767 en recall y
0.3206 en F1. El extremo inferior de precision fue nitrógeno (0.7386); el de
recall y F1 fue potasio (0.6129 y 0.6706). Necrosis letal ocupó el extremo alto
de recall y F1. La desigualdad entre clases es mayor que cualquiera de las
brechas agregadas de calidad de imagen.

Fuente quedó muy cerca: gap 0.3151 en F1 y 0.3599 en recall. Blur produjo 0.0929
de F1; entorno 0.0861; resolución 0.0683. Que entorno tenga un gap global menor
no contradice el sesgo de fondos: sus dos grupos contienen conjuntos de clases
distintos y por eso se necesita la lectura controlada siguiente.

\subsection{Laboratorio frente a campo}

\begin{figure}[H]
  \centering
  \includegraphics[width=0.96\textwidth]{fairness_environment.pdf}
  \caption{A la izquierda, macro métricas globales de grupos con composición de
  clases distinta; a la derecha, recall dentro de las clases presentes en ambos
  entornos.}
  \label{fig:fairness-environment}
\end{figure}

En el resumen global, laboratorio obtuvo F1 0.9352 y campo 0.8491; los recalls
fueron 0.9371 y 0.8376. No se concluye que retirar el fondo cause 0.0861 de
ganancia. Laboratorio no contiene deficiencias y el grupo de campo sí incluye
potasio, la clase más difícil.

Dentro de mancha gris el recall fue 0.9091 en laboratorio y 0.9343 en campo;
para tizón norteño, 0.9023 y 0.9282. Esas dos clases no muestran una caída de
campo. Roya sí: 1.0000 sobre 322 imágenes de laboratorio y 0.3125 sobre solo 16
de campo. El tamaño de ese último grupo impide generalizar el número, pero
es compatible con una brecha de dominio, sin identificar su causa: reconocer roya
sobre fondo controlado no garantiza reconocer sus pocas capturas reales. No se
considera una confirmación causal de dependencia del fondo.

La mitad izquierda no se interpreta como efecto causal del fondo. Laboratorio
solo contiene roya, mancha gris y tizón norteño, mientras campo contiene las
nueve categorías. La comparación dentro de clase reduce ese problema para tres
categorías, aunque todavía puede mezclar fuente, cámara y resolución.

\subsection{Fuente, resolución, nitidez y clases pequeñas}

Entre fuentes, \texttt{maize\_field} tuvo el F1 menor (0.6849), mientras
\texttt{maize\_leaf\_roboflow} alcanzó 1.0000 en su subconjunto. La diferencia
no compara colecciones equivalentes: una fuente puede aportar una sola clase o
un estilo particular. Su utilidad es operativa: antes de añadir nuevas imágenes
de \texttt{maize\_field} conviene revisar por qué casi un tercio de su rendimiento
queda fuera del extremo superior.

El tercil con mayor desenfoque aparente obtuvo F1 0.8022 y el más nítido 0.8951.
La dirección coincide con lo esperado: menos detalle acompaña más errores. En
cambio, resolución no fue monotónica. El tercil medio quedó en 0.8425 y el de
menor resolución en 0.9108. Más píxeles originales no aseguran que la lesión
esté grande, enfocada o bien iluminada; usar resolución como filtro automático
habría descartado casos sin evidencia suficiente.

El indicador de blur es la varianza del Laplaciano de una copia reducida. Una
lesión con textura puede elevarlo y una hoja correctamente enfocada pero lisa
puede reducirlo; por eso se habla de terciles de nitidez aparente. De forma
similar, resolución original no garantiza que la hoja ocupe suficientes
píxeles. Estas variables ayudan a encontrar segmentos débiles, no sustituyen
una auditoría de calidad de captura.

El informe evita la frase «el modelo es justo». Un gap pequeño en este holdout
solo significa que los grupos definidos aquí quedaron cerca. Sin ubicación,
finca, variedad o dispositivo no se puede descartar una brecha regional. El
formato sigue el principio de una ficha de modelo: mostrar soporte, condiciones
y límites junto con la métrica \citep{mitchell2019modelcards}.
